% Curs 11 - Pagina 2 (13-XII-1994).
% Fragment pentru % Curs 11 - Pagina 2 (13-XII-1994).
% Fragment pentru % Curs 11 - Pagina 2 (13-XII-1994).
% Fragment pentru % Curs 11 - Pagina 2 (13-XII-1994).
% Fragment pentru \input{courses/pages/curs11-pagina2-ro}.
% Fara preambul sau comenzi de schimbare fortata a paginii.
% Ordinea si prescurtarile manuscrisului sunt pastrate pe cat posibil.

\subsection*{Curs 11 - Pagina 2}

\noindent Afirm că $|A_0|=1 \Longleftrightarrow
(\forall)\,Q\in A_0$ avem $Q=P$.

\medskip
\[
P\subseteq N_G(Q)=\{a\in G\mid aQa^{-1}=Q\}\supseteq Q.
\]
\[
N_G(Q)\leq G \quad\Longrightarrow\quad |N_G(Q)|\mid p^t m.
\]
\[
|P|=p^t\mid |N_G(Q)|.
\]
\[
\Longrightarrow\quad |N_G(Q)|=p^t m_1,
\qquad p\nmid m_1,\quad m_1\mid m.
\]

\noindent $\Longrightarrow$ $P,Q$ sunt $p$-subgrupuri Sylow ale lui $N_G(Q)$.

\medskip
\noindent Din Sylow II $\Longrightarrow$ $(\exists)\,a\in N_G(Q)$ a.î.
\[
aQa^{-1}=P.
\]
\noindent $a\in N_G(Q) \Longrightarrow aQa^{-1}=Q
\Longrightarrow P=Q \Longrightarrow |A_0|=1$.

\medskip
\noindent\underline{Deducem:}
\[
|\operatorname{Syl}_p(G)|\equiv1\pmod p.
\]
% In manuscris, la aceasta concluzie barele de cardinal sunt omise
% sau neclare: apare "Syl_p(G) = 1 (mod p)".

\bigskip
\noindent\underline{Aplicaţie:}
Fie $G$ un grup finit, $|G|=pq$, $p>q$, $q\nmid p-1$,
$p,q$ prime. \textit{(35)}
\[
\text{Atunci}\qquad G\cong (\mathbb{Z}_{pq},+).
\]
% Cifra/nota din paranteza "(35)" este greu de interpretat in manuscris;
% am pastrat-o ca adnotare, fara a-i atribui o semnificatie.

\noindent $A=\operatorname{Syl}_p(G),\qquad B=\operatorname{Syl}_q(G)$.

\medskip
\noindent Din T. lui Cauchy $\Longrightarrow$ $(\exists)\,a,b\in G$,
$\operatorname{ord}(a)=p$, $\operatorname{ord}(b)=q$.

\[
\begin{aligned}
P&=\langle a\rangle\leq G, & P&\in A,\\
Q&=\langle b\rangle\leq G, & Q&\in B.
\end{aligned}
\]

\noindent Din Sylow III:
$|A|\mid pq$ şi $|B|\mid pq$, deci
\[
\left\{
\begin{aligned}
|A|&\in\{1,p,q,pq\},\\
|B|&\in\{1,p,q,pq\}.
\end{aligned}
\right.
\]

\noindent Arăt că $|A|=|B|=1$.

\medskip
\noindent Consider $[a,b]=aba^{-1}b^{-1}$.

\medskip
\noindent $P\trianglelefteq G$,\quad
$xPx^{-1}\in\operatorname{Syl}_p(G)$, $(\forall)\,x\in G$.
% Pagina se opreste aici. Manuscrisul nu cuprinde inca demonstrarea
% afirmatiei |A|=|B|=1 sau concluzia aplicatiei.
.
% Fara preambul sau comenzi de schimbare fortata a paginii.
% Ordinea si prescurtarile manuscrisului sunt pastrate pe cat posibil.

\subsection*{Curs 11 - Pagina 2}

\noindent Afirm că $|A_0|=1 \Longleftrightarrow
(\forall)\,Q\in A_0$ avem $Q=P$.

\medskip
\[
P\subseteq N_G(Q)=\{a\in G\mid aQa^{-1}=Q\}\supseteq Q.
\]
\[
N_G(Q)\leq G \quad\Longrightarrow\quad |N_G(Q)|\mid p^t m.
\]
\[
|P|=p^t\mid |N_G(Q)|.
\]
\[
\Longrightarrow\quad |N_G(Q)|=p^t m_1,
\qquad p\nmid m_1,\quad m_1\mid m.
\]

\noindent $\Longrightarrow$ $P,Q$ sunt $p$-subgrupuri Sylow ale lui $N_G(Q)$.

\medskip
\noindent Din Sylow II $\Longrightarrow$ $(\exists)\,a\in N_G(Q)$ a.î.
\[
aQa^{-1}=P.
\]
\noindent $a\in N_G(Q) \Longrightarrow aQa^{-1}=Q
\Longrightarrow P=Q \Longrightarrow |A_0|=1$.

\medskip
\noindent\underline{Deducem:}
\[
|\operatorname{Syl}_p(G)|\equiv1\pmod p.
\]
% In manuscris, la aceasta concluzie barele de cardinal sunt omise
% sau neclare: apare "Syl_p(G) = 1 (mod p)".

\bigskip
\noindent\underline{Aplicaţie:}
Fie $G$ un grup finit, $|G|=pq$, $p>q$, $q\nmid p-1$,
$p,q$ prime. \textit{(35)}
\[
\text{Atunci}\qquad G\cong (\mathbb{Z}_{pq},+).
\]
% Cifra/nota din paranteza "(35)" este greu de interpretat in manuscris;
% am pastrat-o ca adnotare, fara a-i atribui o semnificatie.

\noindent $A=\operatorname{Syl}_p(G),\qquad B=\operatorname{Syl}_q(G)$.

\medskip
\noindent Din T. lui Cauchy $\Longrightarrow$ $(\exists)\,a,b\in G$,
$\operatorname{ord}(a)=p$, $\operatorname{ord}(b)=q$.

\[
\begin{aligned}
P&=\langle a\rangle\leq G, & P&\in A,\\
Q&=\langle b\rangle\leq G, & Q&\in B.
\end{aligned}
\]

\noindent Din Sylow III:
$|A|\mid pq$ şi $|B|\mid pq$, deci
\[
\left\{
\begin{aligned}
|A|&\in\{1,p,q,pq\},\\
|B|&\in\{1,p,q,pq\}.
\end{aligned}
\right.
\]

\noindent Arăt că $|A|=|B|=1$.

\medskip
\noindent Consider $[a,b]=aba^{-1}b^{-1}$.

\medskip
\noindent $P\trianglelefteq G$,\quad
$xPx^{-1}\in\operatorname{Syl}_p(G)$, $(\forall)\,x\in G$.
% Pagina se opreste aici. Manuscrisul nu cuprinde inca demonstrarea
% afirmatiei |A|=|B|=1 sau concluzia aplicatiei.
.
% Fara preambul sau comenzi de schimbare fortata a paginii.
% Ordinea si prescurtarile manuscrisului sunt pastrate pe cat posibil.

\subsection*{Curs 11 - Pagina 2}

\noindent Afirm că $|A_0|=1 \Longleftrightarrow
(\forall)\,Q\in A_0$ avem $Q=P$.

\medskip
\[
P\subseteq N_G(Q)=\{a\in G\mid aQa^{-1}=Q\}\supseteq Q.
\]
\[
N_G(Q)\leq G \quad\Longrightarrow\quad |N_G(Q)|\mid p^t m.
\]
\[
|P|=p^t\mid |N_G(Q)|.
\]
\[
\Longrightarrow\quad |N_G(Q)|=p^t m_1,
\qquad p\nmid m_1,\quad m_1\mid m.
\]

\noindent $\Longrightarrow$ $P,Q$ sunt $p$-subgrupuri Sylow ale lui $N_G(Q)$.

\medskip
\noindent Din Sylow II $\Longrightarrow$ $(\exists)\,a\in N_G(Q)$ a.î.
\[
aQa^{-1}=P.
\]
\noindent $a\in N_G(Q) \Longrightarrow aQa^{-1}=Q
\Longrightarrow P=Q \Longrightarrow |A_0|=1$.

\medskip
\noindent\underline{Deducem:}
\[
|\operatorname{Syl}_p(G)|\equiv1\pmod p.
\]
% In manuscris, la aceasta concluzie barele de cardinal sunt omise
% sau neclare: apare "Syl_p(G) = 1 (mod p)".

\bigskip
\noindent\underline{Aplicaţie:}
Fie $G$ un grup finit, $|G|=pq$, $p>q$, $q\nmid p-1$,
$p,q$ prime. \textit{(35)}
\[
\text{Atunci}\qquad G\cong (\mathbb{Z}_{pq},+).
\]
% Cifra/nota din paranteza "(35)" este greu de interpretat in manuscris;
% am pastrat-o ca adnotare, fara a-i atribui o semnificatie.

\noindent $A=\operatorname{Syl}_p(G),\qquad B=\operatorname{Syl}_q(G)$.

\medskip
\noindent Din T. lui Cauchy $\Longrightarrow$ $(\exists)\,a,b\in G$,
$\operatorname{ord}(a)=p$, $\operatorname{ord}(b)=q$.

\[
\begin{aligned}
P&=\langle a\rangle\leq G, & P&\in A,\\
Q&=\langle b\rangle\leq G, & Q&\in B.
\end{aligned}
\]

\noindent Din Sylow III:
$|A|\mid pq$ şi $|B|\mid pq$, deci
\[
\left\{
\begin{aligned}
|A|&\in\{1,p,q,pq\},\\
|B|&\in\{1,p,q,pq\}.
\end{aligned}
\right.
\]

\noindent Arăt că $|A|=|B|=1$.

\medskip
\noindent Consider $[a,b]=aba^{-1}b^{-1}$.

\medskip
\noindent $P\trianglelefteq G$,\quad
$xPx^{-1}\in\operatorname{Syl}_p(G)$, $(\forall)\,x\in G$.
% Pagina se opreste aici. Manuscrisul nu cuprinde inca demonstrarea
% afirmatiei |A|=|B|=1 sau concluzia aplicatiei.
.
% Fara preambul sau comenzi de schimbare fortata a paginii.
% Ordinea si prescurtarile manuscrisului sunt pastrate pe cat posibil.

\subsection*{Curs 11 - Pagina 2}

\noindent Afirm că $|A_0|=1 \Longleftrightarrow
(\forall)\,Q\in A_0$ avem $Q=P$.

\medskip
\[
P\subseteq N_G(Q)=\{a\in G\mid aQa^{-1}=Q\}\supseteq Q.
\]
\[
N_G(Q)\leq G \quad\Longrightarrow\quad |N_G(Q)|\mid p^t m.
\]
\[
|P|=p^t\mid |N_G(Q)|.
\]
\[
\Longrightarrow\quad |N_G(Q)|=p^t m_1,
\qquad p\nmid m_1,\quad m_1\mid m.
\]

\noindent $\Longrightarrow$ $P,Q$ sunt $p$-subgrupuri Sylow ale lui $N_G(Q)$.

\medskip
\noindent Din Sylow II $\Longrightarrow$ $(\exists)\,a\in N_G(Q)$ a.î.
\[
aQa^{-1}=P.
\]
\noindent $a\in N_G(Q) \Longrightarrow aQa^{-1}=Q
\Longrightarrow P=Q \Longrightarrow |A_0|=1$.

\medskip
\noindent\underline{Deducem:}
\[
|\operatorname{Syl}_p(G)|\equiv1\pmod p.
\]
% In manuscris, la aceasta concluzie barele de cardinal sunt omise
% sau neclare: apare "Syl_p(G) = 1 (mod p)".

\bigskip
\noindent\underline{Aplicaţie:}
Fie $G$ un grup finit, $|G|=pq$, $p>q$, $q\nmid p-1$,
$p,q$ prime. \textit{(35)}
\[
\text{Atunci}\qquad G\cong (\mathbb{Z}_{pq},+).
\]
% Cifra/nota din paranteza "(35)" este greu de interpretat in manuscris;
% am pastrat-o ca adnotare, fara a-i atribui o semnificatie.

\noindent $A=\operatorname{Syl}_p(G),\qquad B=\operatorname{Syl}_q(G)$.

\medskip
\noindent Din T. lui Cauchy $\Longrightarrow$ $(\exists)\,a,b\in G$,
$\operatorname{ord}(a)=p$, $\operatorname{ord}(b)=q$.

\[
\begin{aligned}
P&=\langle a\rangle\leq G, & P&\in A,\\
Q&=\langle b\rangle\leq G, & Q&\in B.
\end{aligned}
\]

\noindent Din Sylow III:
$|A|\mid pq$ şi $|B|\mid pq$, deci
\[
\left\{
\begin{aligned}
|A|&\in\{1,p,q,pq\},\\
|B|&\in\{1,p,q,pq\}.
\end{aligned}
\right.
\]

\noindent Arăt că $|A|=|B|=1$.

\medskip
\noindent Consider $[a,b]=aba^{-1}b^{-1}$.

\medskip
\noindent $P\trianglelefteq G$,\quad
$xPx^{-1}\in\operatorname{Syl}_p(G)$, $(\forall)\,x\in G$.
% Pagina se opreste aici. Manuscrisul nu cuprinde inca demonstrarea
% afirmatiei |A|=|B|=1 sau concluzia aplicatiei.
